\section{অধ্যায়ের সারসংক্ষেপ ও গুরুত্বপূর্ণ পরিভাষা}

\par\smallskip\noindent
\begin{minipage}[t]{.46\linewidth}\vspace{0pt}\small
\subsection{গুরুত্বপূর্ণ পরিভাষা}
\begin{tabularx}{\linewidth}{|p{20mm}|Y|}\hline
\chSixHeader পরিভাষা & এক বাক্যে মনে রাখি\\\hline


Database / DBMS & সম্পর্কিত ডেটার সংগ্রহ / তা পরিচালনার সফটওয়্যার\\ \hline
Field / Record & বৈশিষ্ট্যের কলাম / একটি এন্টিটির মানসমূহের সারি\\ \hline
Primary Key / Foreign Key & অনন্য শনাক্তকারী / সম্পর্কিত কী-এর রেফারেন্স\\ \hline
NULL & মান অনুপস্থিত বা অজানা; ০ নয়\\ \hline
Validation & গ্রহণযোগ্য মানের নিয়ম প্রয়োগ\\ \hline
Query / SQL & ডেটাবেজে অনুরোধ / নির্দেশ প্রকাশের ভাষা\\ \hline
Sorting / Indexing & ফল সাজানো / অনুসন্ধানের সহায়ক কাঠামো\\ \hline
Transaction / ACID & যৌক্তিক লেনদেন / নির্ভরযোগ্যতার চার বৈশিষ্ট্য\\ \hline
NoSQL / Big Data & বিকল্প ডেটা মডেল / বৃহৎ ও জটিল তথ্য সামলানোর প্রসঙ্গ\\ \hline
Backup / Recovery & অনুলিপি রাখা / কার্যকর অবস্থায় ফেরানো\\ \hline
Encryption & কী দিয়ে মূল তথ্যকে গোপন সংকেতে রূপান্তর\\ \hline

\end{tabularx}
\end{minipage}\hfill
\begin{minipage}[t]{.51\linewidth}\vspace{0pt}\small
\subsection{SQL নির্দেশ এক নজরে}
\begin{tabularx}{\linewidth}{|p{27mm}|X|}\hline
নির্দেশ & কাজ\\ \hline
CREATE / ALTER / DROP & কাঠামো তৈরি / পরিবর্তন / অপসারণ\\ \hline
SELECT / WHERE & ফিল্ড নির্বাচন / রেকর্ডের শর্ত\\ \hline
INSERT / UPDATE / DELETE & রেকর্ড যোগ / মান সংশোধন / রেকর্ড বাদ\\ \hline
GROUP BY / HAVING & দল তৈরি / দলের শর্ত\\ \hline
ORDER BY / CREATE INDEX & ফল সাজানো / অনুসন্ধানের সহায়ক কাঠামো তৈরি\\ \hline
COMMIT / ROLLBACK & লেনদেন স্থায়ী / চলমান পরিবর্তন বাতিল\\ \hline
\end{tabularx}
\end{minipage}\par\smallskip

\subsection{MCQ-তে ভুল এড়ানোর তথ্য}
SQL = Structured Query Language; SELECT রেকর্ড নির্বাচন করে।

\textbf{কী ও ডেটা টাইপ:} প্রাইমারি কী একাধিক ফিল্ডের হতে পারে; সব কী AutoNumber নয়। ফরেন কী-তে ডুপ্লিকেট হতে পারে; সব ফরেন কী অবশ্যই NULL হবে—এ কথাও ভুল। Date/Time তারিখের জন্য, Currency অর্থের জন্য, Text পরিচিতি/মোবাইলের জন্য উপযোগী।

\textbf{কুয়েরি ও ইনডেক্স:} WHERE শর্ত দেয়, ORDER BY সাজায়; ইনডেক্স ফলের ক্রম নিশ্চিত করে না। COUNT(*) সারি গোনে; COUNT(field) NULL বাদ দেয়। DELETE রেকর্ড বাদ দেয়; DROP TABLE টেবিল বাদ দেয়।

\textbf{পুনরাবৃত্তি ও অনন্যতা:} একই মান দুবার দেখা মানেই নকশা ভুল নয়; কোন ফিল্ড বা সমন্বয় অনন্য হওয়া দরকার সেটি দেখো।

Student(StudentID, Name) ও Result(ResultID, StudentID, Marks) থেকে বেশি নম্বর পাওয়া শিক্ষার্থীর নাম দেখাতে কী কী ধারণা লাগবে? \textbf{ইঙ্গিত:} কী, 1:N, JOIN, WHERE, ORDER BY।
